\documentclass[10pt, a4paper]{article}

\usepackage{setspace}
\setstretch{1.2}      % 줄간격 약간 축소
\usepackage{fontspec}
\usepackage{geometry}
% 여백 유지
\geometry{left=1.6cm, right=1.6cm, top=1.6cm, bottom=1.6cm}
\usepackage{titlesec}
\usepackage{enumitem}
\usepackage{hyperref}
\usepackage{kotex}
\usepackage{array}
\usepackage{xcolor}
\usepackage{fontawesome}

% 폰트 설정 (시스템에 Roboto 폰트가 설치되어 있어야 합니다)
\setmainfont{Roboto}[
    UprightFont = *-Regular,
    BoldFont = *-Bold,
    ItalicFont = *-Italic
]
\setsansfont{Roboto}

% [핵심] 얇고 작은 폰트를 위한 별도 정의 (Roboto Light 사용)
\newfontfamily\lightfont{Roboto Light}

% [핵심] 설명글용 스타일 매크로 정의 (8pt 크기, 얇은 폰트, 줄간격 10pt)
\newcommand{\descstyle}{\lightfont\fontsize{9.5pt}{10pt}\selectfont}

% 리스트 설정 (전역)
\setlist{leftmargin=1em, itemsep=0.3em} 

% 색상 정의
\definecolor{bonusSteelBlue}{RGB}{70, 130, 180}
\definecolor{darkGray}{RGB}{60, 60, 60}
\definecolor{lightGray}{RGB}{150, 150, 150}

% 하이퍼링크 설정
\hypersetup{
    colorlinks=true,
    linkcolor=bonusSteelBlue,
    urlcolor=bonusSteelBlue,
}

% 섹션 스타일
\titleformat{\section}
{\large\bfseries\color{darkGray}}
{}
{0em}
{}[\titlerule]

% 커스텀 명령어
\newcommand{\jobtitle}[2]{{\normalsize\textbf{#1}} \hfill {\descstyle\textit{#2}}}
\newcommand{\dates}[1]{\hfill {\descstyle\color{lightGray} #1}}

\begin{document}
\pagestyle{empty}

% --- HEADER START ---
\noindent
\begin{minipage}[t]{0.61 \textwidth} 
    {\huge\bfseries 두혁진} \\[0.5em]
    {\Large\textcolor{bonusSteelBlue}{system software engineer}} 
    \vspace{0.8em}
    
    {\descstyle 
    시스템의 기저 원리를 파고들며 운영 업무를 자동화하여 안정적인 서비스를 만드는 엔지니어입니다.
    C/C++ 기반의 견고한 시스템 프로그래밍 역량과 Rust의 메모리 안정성을 활용하여 
    모빌리티 소프트웨어의 신뢰성 확보에 기여하고 싶습니다.}
\end{minipage}%
\hfill
\begin{minipage}[t]{0.28\textwidth} 
    {\huge\bfseries contact} \\
    {\descstyle
    \faPhone \ 010-2805-0547 \\
    \faEnvelope \ \href{mailto:tyeirdoo06@gmail.com}{tyeirdoo06@gmail.com} \\
    \faGithub \ \href{https://github.com/hdoo42}{github.com/hdoo42}
    }
\end{minipage}

\vspace{1.5em}

% --- 본문 시작 ---
\noindent
\begin{minipage}[t]{0.34\textwidth} % 왼쪽 컬럼

    % 1. EDUCATION
    \section*{\huge{education}}
    42 Seoul \\ 
    \dates{2021.11 - 2023.12}
    \vspace{0.3em}
    
    % [수정] '메모리 수동 관리', '엄격한 표준' 강조 -> MISRA C 대응 역량 어필
    {\descstyle
    \begin{itemize}[label=\textbullet]
        \item C std lib, bash, IRC server 등을 직접 재구현하며 CS 기초 체화
        \item 엄격한 코드 규칙 준수 및 메모리 수동 관리를 통해 결함 없는 코드 작성 훈련
        \item 6대 coalition master(학생회장) 역임, 커뮤니티에의 공로 인정받아 IITP원장상 수상
    \end{itemize}
    }


    \vspace{1.5em}

    % 2. SKILLS
    \section*{\huge{skills}}
    Programming \\
    {\descstyle Rust, C, C++, Bash, Python }
    
    \vspace{0.8em}
    
    Spoken Languages \\
    {\descstyle Japanese (Advanced Business) \\
    English (Daily / Simple Business) }     

    \vspace{0.8em}
    System \& OS \\
    {\descstyle Linux (Ubuntu), Shell Scripting }
    
    \vspace{0.8em}
    DevOps \& Tools \\
    {\descstyle Neovim, Git, Ansible, Docker, Prometheus, Grafana }

    \vspace{1.5em}

    % 3. ASK ME ABOUT
    \section*{\huge{ask me about}}
    Linux Kernel Troubleshooting
    \begin{itemize}[label=\textbullet]
        \descstyle 
        \item Issue: Ubuntu 22.04 커널 업데이트 후 하드웨어(밝기 조절) 제어 불가 현상 트러블슈팅
        \item Fix: GRUB 커널 파라미터 옵션 수정을 통해 문제 해결
    \end{itemize}

\end{minipage}
\hfill
\begin{minipage}[t]{0.62\textwidth} % 오른쪽 컬럼

    % 4. WORK EXPERIENCE
    \section*{\huge{work experience}}
    
    \jobtitle{(재)경산이노베이션아카데미}{System Software Engineer} \\
    \dates{2024.02 - 현재}
    \begin{itemize}[label=\textbullet] 
        \descstyle 
        \item Operational Automation (GsRoot)
        \begin{itemize}
            \item Rust/SQLite 기반의 워크스테이션 제어 에이전트 개발 (메모리 점유율 21.2MB 최적화)
            \item 관리자 부재 시간(새벽 등)의 대응 공백을 자동화로 해결하여, 기존 수 분~수 시간이 소요되던 장애 처리를 '즉시'로 단축 및 24시간 무중단 운영(Availability) 달성
        \end{itemize}
        
        \item System Integration (BLINK)
        \begin{itemize}
            \item Biostar2 출입 이벤트를 트리거로 LDAP 계정 상태(Unclose/Close)를 실시간 갱신하는 이벤트 기반 동기화 데몬 개발
        \end{itemize}
        
        \item Infrastructure Monitoring
        \begin{itemize}
            \item Prometheus/Slack 연동을 통해 기존의 수동 점검(Ansible ping)에 의존하던 장애 탐지를 자동화하여, 클러스터 상태 사각지대 제거
            \item 노드 다운 및 디스크 임계치(1GB 미만) 도달 시 alert 설정
        \end{itemize}
        
        \item Technical Documentation
        \begin{itemize}
            \item 사내 기술 지식 공유 및 자산화를 위한 위키(Wiki) 시스템 구축 및 문서화 주도
        \end{itemize}
    \end{itemize}

    \vspace{1.2em}

    \jobtitle{(재)이노베이션 아카데미}{Intern} \\
    \dates{2023.10 - 2023.12}
    \begin{itemize}[label=\textbullet]
        \descstyle 
        \item Ansible을 활용한 클러스터 구성 관리
        \item 교육생 과제 명세 작성
    \end{itemize}

    \vspace{1.0em}

    % 5. KEY PROJECTS
    \section*{\huge{key projects}}
    \vspace{0.2em}
    \textbf{42 Region (Private Cloud Build)}
    \begin{itemize}[label=\textbullet]
        \descstyle 
        \item OpenStack 기반 프라이빗 클라우드 구축 중 네트워크 이슈(L1/L2 VLAN) 트러블슈팅
        \item 가상화 인프라 및 네트워크 아키텍처에 대한 심도 있는 이해
    \end{itemize}

\end{minipage}

\end{document}
